\documentclass{article}
\usepackage[T1]{fontenc}
\usepackage{lmodern}
\usepackage{graphicx}
\usepackage{amsmath,amssymb}
\usepackage[a4paper,margin=2cm]{geometry}
\usepackage[absolute,overlay]{textpos}
\setlength{\TPHorizModule}{1mm}
\setlength{\TPVertModule}{1mm}
\usepackage[indonesian]{babel}
\usepackage{hyperref}
\setlength{\emergencystretch}{2em}
% unit: Halaman 13

\begin{document}

% === Halaman 13 (llm) ===

\begin{itemize}
    \item Jika panjang bit-bit kunci lebih pendek daripada panjang bit-bit pesan, maka bit-bit kunci diulang penggunaannya secara periodik (seperti halnya pada Vigenere Cipher)
    \item Contoh:
    \begin{align*}
    \text{Plainteks} & : 10010010101110101010001110001 \\
    \text{Kunci} & : 11011011011011011011011011011 \\
    \text{Cipherteks} & : 01001001110101110001010101010
    \end{align*}
\end{itemize}

\end{document}
